\scanpage{112}

de $G_T$ sur $\underline F$ équivaut à la donnée d'une opération de $G$ sur $\underline E$ \inserted{commutant à $\mathscr B$}, car $G_T\unsure{-|\longrightarrow}\underline{\mathrm{Aut}}(\underline F)$ équivaut à

\[
G\longrightarrow\prod_{T/S}\underline{\mathrm{Aut}}(\underline F),\quad\text{or}
\]

\[
\prod_{T/S}\underline{\mathrm{Aut}}(\underline F)\simeq
\underline{\mathrm{Aut}}_{\mathscr B}\Bigl(\prod_{T/S}\underline F\Bigr)
=\underline{\mathrm{Aut}}_{\mathscr B}(\underline E).
\]

Donc, la catégorie $C(G_T)$ équivaut à la catégorie des objets de $C(G)$ munis d'une opération de $\mathscr B$, ce qui, si $S$ est affine d'anneau $A$, donc $\unsure{B|\mathscr B|T|\mathcal B}$ défini par $A$-algèbre $B$, équivaut à la \deleted{\unsure{caté|catég.}} donnée d'un objet de $C(G)$ à anneau $B$ d'endomorphismes.

Lorsque $G$ satisfait l'une des conditions a) ou b) (p.ex. $G$ qu. cpct sur un \uline{corps} $k$), \deleted{\unsure{les|la}} catégorie $C(G_T)$ est définie en termes de $C_f(G)$.

De $\overline{\mathrm m}$, si $\varphi_G$ \inserted{$(\varphi_G^f)$} désigne le foncteur oubli de \inserted{$C(G)$} ($C_f(G)$), alors $(C(G_T),\varphi_{G_T})$ est déterminé en termes de $(C_f(G),\varphi_G^f)$.

Et il en est donc de $\overline{\mathrm m}$ de $C_f(G_T),\varphi_{G_T})$, sous-catégorie pleine de $C(G_T)$ formée des $E$ tels que $\varphi_{G_T}(E)$ de d.f. sur $B$.
% archive sheet number: 111
